\chapter{Blackwell's Discounted Decisions on Borel Spaces}
\label{sec:p17-blackwell-general-discounted}

Blackwell's general formulation permits uncountably many states and actions,
arbitrary bounded measurable rewards, and arbitrary measurable transition
probabilities \cite[pp.~226--235]{Blackwell1965Discounted}.
The distinction between a policy good outside a null set for a specified
initial distribution and one good at every initial state is essential.
This chapter supplies the measurable constructions behind that distinction.
The finite model of \cref{sec:p17-sequential-decisions} embeds through
the common-action construction in \cref{thm:p17-sequential-source-bridge},
with the discrete sigma algebras on its finite sets.
No continuity or compactness of
the action optimization is assumed here.

\section{Measurable spaces, factorization and selection}

\begin{hbdefinition}{def:p17-blackwell-borel-model}{The measurable decision model}
A standard Borel space is a Borel subset of a complete separable metric
space, equipped with its trace Borel sigma algebra. All spaces below
are nonempty. Products carry their coordinate product sigma algebras.
For these spaces they agree with the Borel sigma algebras of the
product topologies. Write \(M(X)\) for bounded real Borel functions on \(X\),
\(P(X)\) for probability measures on \(X\), and \(Q(Y\mid X)\) for
probability kernels from \(X\) to \(Y\).
For \(k\in Q(Y\mid X)\) and \(v\in M(X\times Y)\), put
\[
 (kv)(x)=\int_Y v(x,y)k(dy\mid x).
\]
Blackwell's term \emph{Baire function} refers here to Borel measurability:
on a metric space the sigma algebra generated by continuous real
functions is its Borel sigma algebra, since open sets are countable
unions of positive distance level sets of their closed complements.

The decision data are standard Borel \(S,A\), a kernel
\(q\in Q(S\mid S\times A)\), a reward \(r\in M(S\times A\times S)\),
and \(0\leq\beta<1\). Put \(R=\|r\|_\infty\).
On day \(n\geq1\), the observed history is
\[
 h_n=(s_1,a_1,\ldots,a_{n-1},s_n)\in
 H_n=S\times(A\times S)^{n-1}.
\]
A policy is a sequence \(\pi_n\in Q(A\mid H_n)\).
After sampling \(a_n\), sample \(s_{n+1}\) from
\(q(\cdot\mid s_n,a_n)\), and receive \(r(s_n,a_n,s_{n+1})\).
Its date-one value uses discount \(\beta^{n-1}\).
All kernels are specified also at histories having probability zero.
A deterministic Markov policy is \((f_1,f_2,\ldots)\) with Borel
\(f_n:S\to A\); the stationary policy \(f^\infty\) uses one such map
at every date. This is the source's nonrandomized meaning of Markov.
\end{hbdefinition}

\begin{hbtheorem}{thm:p17-blackwell-factorization}{proved=annex:p17.blackwell.factorization}{Factorization and parameter-dependent infinite products}
Every \(m\in P(X\times Y)\) on standard Borel spaces factors as
\(m(dx,dy)=p(dx)k(dy\mid x)\). Its first marginal \(p\) is unique;
the kernel is unique outside one \(p\)-null set, as a measure on
all Borel target events simultaneously.
Every bounded Borel integrand has the iterated integral specified above.

For standard Borel coordinate spaces, an initial probability and
successive Borel kernels, there is a unique countable product law
with the prescribed finite iterated distributions. If the initial
probability and these kernels depend measurably on a parameter \(z\),
the resulting laws form a probability kernel in \(z\).
Consequently each policy has a future kernel
\(K^\pi\in Q((A\times S)^{\mathbb N_+}\mid S)\), and
\[
 I^\pi(s)=E_s^\pi\sum_{n=1}^{\infty}
       \beta^{n-1}r(S_n,A_n,S_{n+1})\in M(S).
\]
One has \(\|I^\pi\|_\infty\leq R/(1-\beta)\).
Truncation after \(N\) rewards has error at most
\(R\beta^N/(1-\beta)\). Policies agreeing through day \(N\)
have value difference at most \(2R\beta^N/(1-\beta)\).
These assertions also hold jointly in any additional Borel
parameter retained in the policy kernels.
\end{hbtheorem}

\begin{hblemma}{thm:p17-blackwell-positive-selection}{proved=annex:p17.blackwell.positive.selection}{A Borel selector from positive kernel sections}
Let \(D\subseteq X\times Y\) be Borel and
\(k\in Q(Y\mid X)\). If \(k(D_x\mid x)>0\) for every \(x\),
there is a Borel \(f:X\to Y\) with \((x,f(x))\in D\) for
every \(x\). The proof constructs a compatible complete metric
on the graph space, keeping its Borel sets unchanged, and selects
through nested balls of positive conditional probability.
It requires the positive-section kernel hypothesis; nonempty
sections alone do not supply this construction.
\end{hblemma}

\begin{hblemma}{thm:p17-blackwell-selection}{proved=annex:p17.blackwell.selection}{The exact two conclusions of Blackwell's selection lemma}
For \(k\in Q(Y\mid X)\), \(v\in M(X\times Y)\), and
\(\eta>0\), there is a Borel \(f:X\to Y\) satisfying, for every \(x\),
\[
 v(x,f(x))\geq(kv)(x),\qquad
 k(\{y:v(x,y)\geq v(x,f(x))+\eta\}\mid x)=0.
\]
Both conclusions are pointwise in the conditioning argument.
The second concerns the distribution specified by \(k\), rather than
all actions in \(Y\). A degenerate kernel is equivalently
\(x\mapsto\delta_{f(x)}\) for a Borel \(f\).
\end{hblemma}

\section{What an initial distribution permits}

\begin{hbdefinition}{def:p17-blackwell-optimality}{Three quantifiers over initial states}
For \(\rho\in P(S)\) and \(\varepsilon\geq0\), a policy \(\pi^*\) is
\((\rho,\varepsilon)\)-optimal when
\[
 \forall\pi\quad
 \rho\{s:I^\pi(s)>I^{\pi^*}(s)+\varepsilon\}=0.
\]
The exceptional set can depend on the competing policy.
It is \(\varepsilon\)-optimal when the inequality
\(I^\pi(s)\leq I^{\pi^*}(s)+\varepsilon\) holds for all
\(\pi,s\). It is optimal when this holds with \(\varepsilon=0\).
A policy \((\rho,\varepsilon)\)-dominates a specified \(\pi\)
when \(I^{\pi^*}\geq I^\pi-\varepsilon\) \(\rho\)-almost everywhere.
The policy produced for one \(\rho\) need not work for another.
Uniform \(\varepsilon\)-optimality is equivalent to
\((\rho,\varepsilon)\)-optimality for every \(\rho\):
the converse follows by taking point masses at individual states.
A fixed policy is \(\rho\)-optimal if it is
\((\rho,\varepsilon)\)-optimal for every \(\varepsilon>0\).
For each competitor this is equivalent to zero-tolerance domination,
by intersecting the good sets for \(\varepsilon=1/n\).
\end{hbdefinition}

\begin{hbtheorem}{thm:p17-blackwell-prior-optimal}{proved=annex:p17.blackwell.prior.optimal}{Blackwell's Theorem 1: approximate optimality for a prior}
For every \(\rho\in P(S)\) and \(\varepsilon>0\), a
\((\rho,\varepsilon)\)-optimal policy exists.
It is obtained by choosing policies whose integrated values approach
\(\sup_\pi\int I^\pi\,d\rho\), taking their countable pointwise
supremum, and using a Borel partition of initial states to choose
a policy within \(\varepsilon\) of that supremum.
This is an existence construction and does not promise an effective
search for the near maximizing policies from an arbitrary Borel input.
\end{hbtheorem}

\begin{hbtheorem}{thm:p17-blackwell-prior-markov}{proved=annex:p17.blackwell.prior.markov}{Blackwell's Theorem 2: removing history for a fixed prior}
For every policy \(\pi\), initial distribution \(\rho\), and
\(\varepsilon>0\), there is a deterministic Markov policy
\(\widehat\pi\) with
\[
             I^{\widehat\pi}\geq I^\pi-\varepsilon
                       \quad\rho\text{-almost everywhere}.
\]
In particular a \((\rho,\varepsilon)\)-optimal deterministic
Markov policy exists. The construction truncates the history-dependent
part with the signed tail bound, then replaces its decisions backwards.
At each replacement it uses the conditional distribution of the
current action given the current state under the specified prior
and the current prefix policy.
\end{hbtheorem}

\begin{hbtheorem}{thm:p17-blackwell-projection-boundaries}{proved=annex:p17.blackwell.projection.boundaries}{Three exact failures of stronger existence claims}
There exists a Borel \(B\subseteq[0,1]^2\) whose projection
\(D=\{x:\exists y\ (x,y)\in B\}\) is not Borel. Such a set
can be constructed by the universal closed-set diagonal argument
given in the proof. It yields the following original examples.
\begin{enumerate}
\item With one state, actions \(a=1,2,\ldots\), and reward
      \(1-1/a\), the value supremum is \(1/(1-\beta)\), and
      no randomized history policy attains it.
\item With \(S=A=[0,1]\), a fixed state, and reward
      \(\mathbf1_B(s,a)\), no policy is \(\varepsilon\)-optimal
      for \(\varepsilon<1\).
\item With the tagged disjoint state space
      \(S=(\{0\}\times B)\,\sqcup\,(\{1\}\times[0,1])\),
      let \((0,x,y)\) move to \((1,x)\) and let \((1,x)\)
      stay there. The action space is \([0,1]\);
      the first transition earns zero, and transitions from
      \((1,x)\) earn \(\mathbf1_B(x,a)\).
      A policy remembering the initial \(y\) earns
      \(\beta/(1-\beta)\) from every \((0,x,y)\).
      For \(0<\beta<1\), no Markov policy, even a randomized
      one, dominates this policy within any \(\varepsilon<\beta\)
      at all initial states.
\end{enumerate}
The third assertion concerns positive discount. At \(\beta=0\),
the initial transition reward in that example is zero under every policy.
\end{hbtheorem}

\section{Generated choices and the full general theorem}

\begin{hbdefinition}{def:p17-blackwell-generated}{Stationary and generated operators}
For Borel \(f:S\to A\), define
\[
 (T_f u)(s)=\int_S[r(s,f(s),s')+\beta u(s')]q(ds'\mid s,f(s)).
\]
For a constant action \(a\), write \(T_a\).
Fix a deterministic Markov sequence \(\pi=(f_1,f_2,\ldots)\).
A map \(g\) is \(\pi\)-generated if a countable Borel partition
\((E_n)\) satisfies \(g=f_n\) on \(E_n\).
A deterministic Markov policy is generated if every one of its maps is.
Put
\[
                (Uu)(s)=\sup_{n\geq1}(T_{f_n}u)(s).
\]
The supremum is countable, so \(Uu\) is Borel.
The associated countable-action model has action \(n\in\mathbb N_+\),
transition \(\widehat q(\cdot\mid s,n)=q(\cdot\mid s,f_n(s))\),
and reward \(\widehat r(s,n,s')=r(s,f_n(s),s')\).
The notation \((g,\pi)\) means use \(g\) on the first day and
then restart \(\pi\) at the new state, with its own history clock.
The earlier first state and action are erased in this restart.
\end{hbdefinition}

\begin{hbtheorem}{thm:p17-blackwell-generated-operators}{proved=annex:p17.blackwell.generated.operators}{Blackwell's Theorems 3--5 on the full bounded Borel space}
The operators \(T_f,U\) preserve \(M(S)\), are monotone,
and satisfy \(B(u+c)=Bu+\beta c\) for real constants \(c\),
where \(B\) denotes either operator in this assertion.
They are contractions with modulus \(\beta\).
The fixed point of \(T_f\) is \(I^{f^\infty}\);
\(U\) has a unique fixed point \(u^*\in M(S)\).
For every \(u\in M(S)\), every \(\eta>0\), and every generated \(g\),
\[
 T_g u\leq Uu,\qquad
 \text{there is a generated }f\text{ with }T_fu\geq Uu-\eta .
\]
For every \(N\geq0\), \(U^Nu\) is the supremum of the
\(N\)-day returns with terminal reward \(u\) over generated policies.
For every positive tolerance there is one generated policy achieving
that tolerance simultaneously at every state.
Also \(T_gI^\pi=I^{(g,\pi)}\), including restarted history policies.
The deterministic Markov policies of the associated countable-action
model correspond, with identical state laws and rewards, to generated
policies in the original model; different indices can describe the
same original action.
Successive approximation satisfies
\[
 \|U^n u-u^*\|_\infty\leq\beta^n\|u-u^*\|_\infty,\qquad
 \|U^n u-u^*\|_\infty
       \leq\frac{\beta^n}{1-\beta}\|Uu-u\|_\infty .
\]
\end{hbtheorem}

\begin{hbtheorem}{thm:p17-blackwell-generated-value}{proved=annex:p17.blackwell.generated.value}{Theorem 6(a): the value of all generated plans}
For the preceding sequence and fixed point,
\[
 I^{\pi'}\leq u^*\quad\text{for every generated policy }\pi'.
\]
For every \(\varepsilon>0\) there is a generated stationary
\(f^\infty\) with \(I^{f^\infty}\geq u^*-\varepsilon\).
Every generated \(f\) satisfying
\(T_fu^*\geq u^*-\varepsilon(1-\beta)\) has this property.
\end{hbtheorem}

\begin{hbtheorem}{thm:p17-blackwell-stationary-prior}{proved=annex:p17.blackwell.stationary.prior}{Theorem 6(b)--(c): stationary replacement with the exact tolerances}
For every \(\rho\in P(S)\) and \(\varepsilon>0\), there is a
\((\rho,\varepsilon)\)-optimal deterministic stationary policy.
If an \(\varepsilon\)-optimal policy exists for some
\(\varepsilon\geq0\), there is an
\(\varepsilon/(1-\beta)\)-optimal deterministic stationary policy.
In particular the existence of any optimal policy implies the
existence of a deterministic stationary optimal policy.
The \(\varepsilon=0\) conclusion uses the exact mean improvement
in the selection lemma; it does not require a maximum over \(A\).
\end{hbtheorem}

\begin{hbtheorem}{thm:p17-blackwell-general-verification}{proved=annex:p17.blackwell.general.verification}{Theorem 6(d)--(f): certificates and the optimality equation}
If \(u\in M(S)\) satisfies \(T_a u\leq u\) for every \(a\in A\),
then \(I^\pi\leq u\) for every policy \(\pi\).
If an \(\varepsilon\)-optimal policy exists for every
\(\varepsilon>0\), then
\[
 V(s)=\sup_\pi I^\pi(s)
\]
is bounded Borel and satisfies \(V(s)=\sup_{a\in A}(T_aV)(s)\).
A policy is optimal if and only if its own return satisfies this
equation. Among bounded Borel functions there is at most one
solution of the optimality equation, even when the pointwise
supremum operator fails to preserve Borel functions in general.
This uniqueness statement compares two existing solutions and
does not assert that an arbitrary uncountable action supremum
is measurable.
\end{hbtheorem}

\section{Essentially countable choices and improvement}

\begin{hbdefinition}{def:p17-blackwell-essential-actions}{Two equivalence notions and the source's covering conditions}
Actions \(a,b\) are strongly equivalent at state \(s\) when
\[
 q(\cdot\mid s,a)=q(\cdot\mid s,b),\qquad
 r(s,a,s')=r(s,b,s')\quad(s'\in S).
\]
This is Blackwell's displayed condition.
They are operator equivalent when \(T_au(s)=T_bu(s)\)
for every \(u\in M(S)\).
Given \(\pi=(f_n)\), the actions are essentially countable
when every \((s,a)\) has some \(n\) strongly equivalent to
\(f_n(s)\) at \(s\).
They are essentially finite when there is a countable Borel
partition \(S=\bigsqcup_{m\geq1}E_m\) such that, for
\(s\in E_m\), every action is strongly equivalent to one of
\(f_1(s),\ldots,f_m(s)\).
The conclusions below also hold with operator equivalence
in place of strong equivalence.
\end{hbdefinition}

\begin{hbtheorem}{thm:p17-blackwell-essential-optimum}{proved=annex:p17.blackwell.essential.optimum}{Theorem 7 on arbitrary Borel state spaces}
Under essential countability, \(Uu=\sup_aT_au\) for every
\(u\in M(S)\). Its fixed point is the value over all randomized
history policies, and a deterministic stationary
\(\varepsilon\)-optimal policy exists for every \(\varepsilon>0\).
Under essential finiteness a deterministic stationary optimal
policy exists. In particular countable and finite common action
sets have the corresponding conclusions on every standard Borel
state space, with arbitrary bounded Borel reward and Borel kernel.

For \(0<\beta<1\), operator equivalence holds exactly when
the transition measures and expected immediate rewards agree.
At \(\beta=0\), only expected immediate rewards need agree.
Thus the stronger pointwise reward condition is sufficient
but is not necessary for equality of the action operators.
\end{hbtheorem}

\begin{hbtheorem}{thm:p17-blackwell-general-improvement}{proved=annex:p17.blackwell.general.improvement}{Theorem 8: improvement on the original general domain}
Let \(g:S\to A\) be Borel and let \(\pi\) be a policy.
If \(I^{(g,\pi)}\geq I^\pi\) pointwise, then
\[
              I^{g^\infty}\geq I^{(g,\pi)}\geq I^\pi.
\]
For Borel \(f,g:S\to A\), define \(h=f\) where
\(I^{f^\infty}\geq I^{g^\infty}\), and \(h=g\) elsewhere.
Then \(h\) is Borel and
\[
              I^{h^\infty}\geq
                      \max(I^{f^\infty},I^{g^\infty}).
\]
These are pointwise order improvements. Neither assertion
requires a finite action set or an attained global maximizer.
They do not assert finite termination of policy iteration.
\end{hbtheorem}

\begin{hbtheorem}{thm:p17-blackwell-borel-examples}{proved=annex:p17.blackwell.borel.examples}{Continuous states, infinitely many choices and strict switching}
The following calculations illustrate the exact domains.
The stated rewards specify all realized transitions; give every
impossible transition reward zero to obtain a bounded Borel
function on the full product \(S\times A\times S\).
\begin{enumerate}
\item Let \(S=[0,1]\), let every action leave the state fixed,
      let \(A=\mathbb N_+\), and let \(r(s,n,s)=s-1/n\).
      Then \(V(s)=s/(1-\beta)\), no policy is optimal at any
      state, and the constant action \(n\) is
      \(1/[n(1-\beta)]\)-optimal. From \(u_0=0\), value iteration
      gives \(u_k(s)=s(1-\beta^k)/(1-\beta)\).
\item Let \(S=[0,1]\), \(A=\{0,1\}\), and let \(C\subseteq[0,1]\)
      be any Borel set. Keep the state fixed and reward action
      \(1\) by \(\mathbf1_C(s)\), action \(0\) by
      \(1-\mathbf1_C(s)\). The Borel rule \(f=\mathbf1_C\)
      is optimal, with value \(1/(1-\beta)\), regardless of
      discontinuities at the boundary of \(C\).
\item In a two-state system alternating deterministically between
      \(0,1\), an action matching the current state earns \(1\),
      and the other action earns \(0\). For \(0<\beta<1\),
      the two constant-action policies have values
      \((1,\beta)/(1-\beta^2)\) and
      \((\beta,1)/(1-\beta^2)\). Their switching policy earns
      \(1/(1-\beta)\) at both states, strictly more than their
      pointwise maximum.
\end{enumerate}
\end{hbtheorem}

\begin{proofboundary}{def:p17-blackwell-general-boundary}{What the general discounted method supplies}
The input consists of actual Borel spaces, bounded reward and measurable
probability kernels. The output is a policy existence or verification
statement with its displayed state quantifiers and error tolerance.
Arbitrary Borel input descriptions need not admit numerical integration,
an effective representation of the refined topology, or an algorithm
finding the policies used in the integrated supremum argument.
Effective finite-table computation has the separate method in
\cref{def:p17-sequential-computation}.
Unbounded rewards, undiscounted criteria and universally measurable
policy classes require their own hypotheses and constructions.
The proof supplied here reconstructs the results of the held 1965
article; it does not claim to reproduce the complete separate works
which that article cites.
\end{proofboundary}
